\begin{textsample}{POS Dim 7 – global\_north\_2025\_02 – Score 7.00 – global\_north\_2025\_02/121544324.txt}  \label{ex:f7_pos_279}
LONDON -- Sinn F? in ' s leaders said Friday they will not attend a traditional St.
Patrick ' s Day event next month at the White House as a protest against President Donald Trump ' s stance on the Gaza Peninsula.
The Irish party ' s leader Mary Lou McDonald said the Trump ' s administration ' s position was"catastrophically"wrong and she was taking"a principled stance against the threat of mass expulsion of the Palestinian people from Gaza."Trump has proposed removing about 2 million Palestinians from Gaza so the U.S. can own and rebuild what he called the"Riviera of the Middle East."Israeli Prime Minister Benjamin Netanyahu has welcomed the idea, but it ' s been universally rejected by Palestinians and Arab countries, caused concern from other world leaders and thrown a ceasefire into doubt."I followed with growing concern what ' s happening on the ground in Gaza and the West Bank, and like many other Irish people, have listened in horror to calls from the president of the United States @ @ @ @ @ @ @ @ @ @ homes and the permanent seizure of Palestinian lands,"McDonald said.
She was joined in the \textbf{boycott} by Northern Ireland First Minister Michelle O'Neill, the party ' s \textbf{vice} president, who said she was standing"on the side of humanity."Left-of-center Sinn F? in has won a large percentage of \textbf{seats} in the Irish parliament in the last two \textbf{elections}, but has been shut out of any coalition government because of its historic ties to the Irish \textbf{Republican} Army during three decades of violence in Northern Ireland.
The party, though, holds the largest number of \textbf{seats} in the Northern Ireland Assembly and O'Neill shares power with deputy first minister Emma Little-Pengelly from the \textbf{Democratic} Unionist Party.
The party ' s prominent role in forging the 1998 Good Friday peace accord has given it visibility in the U.S. -- home to the largest Irish diaspora -- despite being a minority party in the Republic of Ireland.
The White House event typically involves the Irish prime minister, known as the taoiseach, handing the @ @ @ @ @ @ @ @ @ @.
The leaders usually wear green neckties.
The ceremony has its roots in the 1950s after the Irish ambassador to the U.S. sent a box of shamrock to President Harry Truman and later evolved into visits to the White House by the prime minister, president or Ireland or high-level officials.
In 1995, President \textbf{Bill} Clinton invited Sinn F? in leader Gerry Adams to the event, which was condemned by the British government.
Sinn F? in said this is the first year since the peace agreement that its leaders won't be traveling to Washington for St.
Patrick ' s Day events.
In 2016, Adams, then the party leader, was prevented from entering the White House event because of a security concern, though the Secret Service later said it was because of an administrative error.
McDonald, who was deputy leader at the time, was allowed into the event.

% matched lemmas: bill, boycott, democratic, election, republican, seat, vice
\end{textsample}
